\section{Proofs}
\label{app:proofs}

The statements below are standard. We include proofs so that the assumptions are explicit.

\paragraph{Proposition~\ref{prop:rrr}.} Write $L_\lambda(W)=\fnorm{(W-M_\lambda)A_\lambda^{1/2}}^2
+\mathrm{const}$ with $M_\lambda=C_\lambda A_\lambda^{-1}$. This is valid when $A_\lambda\succ0$
and follows by completing the square in $\tr(WA_\lambda W^\top)-2\tr(WC_\lambda^\top)$. The
change of variables $N=WA_\lambda^{1/2}$ is a bijection of $\Wr$ onto itself, because
$A_\lambda^{1/2}$ is invertible and preserves rank. Minimising $\fnorm{N-C_\lambda
A_\lambda^{-1/2}}^2$ over $\rank(N)\le r$ is the Eckart--Young problem
\citep{eckart1936approximation}. Its minimisers are the rank-$r$ truncated SVDs, and they are
unique if and only if $\sigma_r>\sigma_{r+1}$. Mapping back gives $W_\lambda$. For
$r\ge\min(p,q)$ no truncation occurs and $W_\lambda=M_\lambda$. \hfill$\square$

\paragraph{Proposition~\ref{prop:unique}.} Let $W^\circ$ minimise $F$ over $\Wr$. Then
$\mathrm{OPT}=F(W^\circ)\ge\sum_k\lambda^\star_kE_k(W^\circ)=L_{\lambda^\star}(W^\circ)\ge
g(\lambda^\star)=\mathrm{OPT}$. Hence $W^\circ$ minimises $L_{\lambda^\star}$, and uniqueness gives
$W^\circ=W_{\lambda^\star}$. \hfill$\square$

\paragraph{Lemma~\ref{lem:pert}.} Let $Y=W-D_k$ with rows $y_i^\top$. Then
$\tr(Y\Delta_kY^\top)=\sum_iy_i^\top\Delta_ky_i$. For symmetric $\Delta_k$,
$|y^\top\Delta_ky|\le\opnorm{\Delta_k}\|y\|^2$. Summing over rows gives
$|\tr(Y\Delta_kY^\top)|\le\opnorm{\Delta_k}\fnorm{Y}^2$, and dividing by $d>0$ proves the claim.
This is the usual spectral-norm bound on a quadratic form. \hfill$\square$

\paragraph{Lemma~\ref{lem:dev}.} For any $W'\in\mathcal W$,
$F(\What)\le\Fhat(\What)+\delta\le\Fhat(W')+\epsilon_{\mathrm{opt}}+\delta\le
F(W')+2\delta+\epsilon_{\mathrm{opt}}$. Taking the infimum over $W'$ gives the claim. This is the
standard uniform-deviation argument for approximate empirical minimisers. \hfill$\square$

\paragraph{Corollary~\ref{cor:fixed}.} Since $|\max_ka_k-\max_kb_k|\le\max_k|a_k-b_k|$,
Lemma~\ref{lem:pert} and $\fnorm{W-D_k}\le B+\fnorm{D_k}$ on $\WrB$ give the bound on
$\sup_{\WrB}|F-\Fhat|$. Next, $\WrB\subseteq\Wr$ implies $\inf_{\Wr}\Fhat\le\inf_{\WrB}\Fhat$,
so the $\epsilon_{\mathrm{opt}}$-optimality of $\What$ over $\Wr$ carries over to $\WrB$ whenever
$\What\in\WrB$. Lemma~\ref{lem:dev} with $\mathcal W=\WrB$ and $\inf_{\WrB}F\le F(W^\star)$ then
completes the proof. The bound applies to our solver's output only a posteriori, when the chosen
$B$ covers both $\What$ and $W^\star$. The solver itself never enforces a norm bound.
\hfill$\square$

\paragraph{Lemma~\ref{lem:denom}.} Write $q=\tr[(W-D_k)S_k(W-D_k)^\top]$ and let $\hat q$ be the
same expression with $\Shat_k$. Then
\[
\frac{q}{d_k}-\frac{\hat q}{\dhat_k}=\frac{q-\hat q}{\dhat_k}+\frac{q}{d_k}\cdot
\frac{\dhat_k-d_k}{\dhat_k}.
\]
By Lemma~\ref{lem:pert}, $|q-\hat q|\le\opnorm{\Delta_k}\fnorm{W-D_k}^2$. Moreover
$|\dhat_k-d_k|=|\tr(D_k\Delta_kD_k^\top)|\le\opnorm{\Delta_k}\fnorm{D_k}^2$. Combining the two
proves the claim. The assumption $\dhat_k\ge d_{\min,k}>0$ is what keeps the bound finite when it
is used uniformly over draws. For a single draw the observed $\dhat_k>0$ suffices. We make no
claim about how large $\opnorm{\Delta_k}$ is with high probability. \hfill$\square$

\section{Implementation and numerical audit}
\label{app:impl}

The implementation is pure Python (standard library only), because no numerical library was
installed in the execution environment.
\begin{itemize}
\item \emph{Inner step.} The inner problem is solved with Proposition~\ref{prop:rrr}, using Jacobi
eigendecompositions and one-sided Jacobi SVDs.
\item \emph{Dual ascent.} It consists of \CfgDualIters{} diminishing-step exponentiated-gradient
steps followed by \CfgPolishIters{} backtracking steps. Every evaluated $g(\lambda)$ is a lower
bound in exact arithmetic.
\item \emph{Primal refinement.} Refinement runs only when the dual gap exceeds $10^{-9}$ in
relative terms. It alternates two convex block problems, one with the column space fixed and one
with the row space fixed, with smoothed local descent on a factorisation. A candidate is accepted
only if it lowers $\Fhat$.
\item \emph{Test suite.} The tests compare the inner solver with a raw-data reduced-rank
regression, a grid search, and multi-start alternating least squares. They check the minimax
bounds against a grid-plus-ellipsoid reference on $2\times q$ problems. They reproduce the
published fair-PCA gap instance \citep[Lemma~6.2]{tantipongpipat2019multicriteria} and the
two-group exactness \citep[Thm.~1.2]{tantipongpipat2019multicriteria}. They also check invariance
under $(B,A)\mapsto(BR,R^{-1}A)$.
\end{itemize}
Passing tests are evidence about the implementation on the tested instances, not proofs.

\section{Additional results}
\label{app:results}

\begin{table}[h]
\caption{Population worst-task error $F$ for every stage-2 cell, one row per teacher and
calibration draw. ``orig.\ (seen)'' is the stage-1 draw, whose test results were seen before the
diagnostic was designed.}
\label{tab:perdraw}
\centering\small
\resizebox{\linewidth}{!}{% AUTO-GENERATED
\begin{tabular}{llcccccc}
\toprule
Teacher & Draw & \multicolumn{2}{c}{Emp/Emp} & \multicolumn{2}{c}{Emp/Pop} & \multicolumn{2}{c}{Pop/Emp} \\
 & & Unif. & Minimax & Unif. & Minimax & Unif. & Minimax \\
\midrule
0 & orig.\ (seen) & 0.657 & 0.726 & 0.765 & 0.712 & 0.609 & 0.600 \\
0 & r101 & 0.635 & 0.550 & 0.739 & 0.665 & 0.632 & 0.586 \\
0 & r102 & 0.651 & 0.645 & 0.628 & 0.686 & 0.581 & 0.533 \\
0 & r103 & 0.584 & 0.550 & 0.638 & 0.754 & 0.754 & 0.727 \\
0 & Pop/Pop (oracle) & \multicolumn{6}{c}{Unif.\ 0.557, Minimax 0.491} \\
\midrule
1 & orig.\ (seen) & 0.655 & 0.768 & 0.636 & 0.813 & 0.719 & 0.669 \\
1 & r101 & 0.616 & 0.620 & 0.634 & 0.684 & 0.685 & 0.642 \\
1 & r102 & 0.699 & 0.619 & 0.751 & 0.647 & 0.634 & 0.681 \\
1 & r103 & 0.660 & 0.693 & 0.671 & 0.717 & 0.646 & 0.566 \\
1 & Pop/Pop (oracle) & \multicolumn{6}{c}{Unif.\ 0.658, Minimax 0.520} \\
\midrule
2 & orig.\ (seen) & 0.780 & 1.081 & 0.736 & 0.909 & 0.640 & 0.757 \\
2 & r101 & 0.694 & 0.821 & 0.790 & 0.816 & 0.778 & 0.822 \\
2 & r102 & 0.755 & 0.726 & 0.676 & 0.760 & 0.742 & 0.719 \\
2 & r103 & 0.694 & 0.822 & 0.761 & 0.908 & 0.701 & 0.648 \\
2 & Pop/Pop (oracle) & \multicolumn{6}{c}{Unif.\ 0.673, Minimax 0.597} \\
\bottomrule
\end{tabular}
}
\end{table}

\begin{table}[h]
\caption{Stage-1 synthetic fixture: all methods at rank $r=\CfgR$, means over \CfgNTeachers{}
teachers, evaluated on the test sample. $\dagger$: unofficial simplified re-implementation.
$\ddagger$: exceeds the rank budget and is shown for reference only. $\S$: factor averaging
depends on the factorisation gauge and is a diagnostic only. Weight-space baselines were tuned on
the development split, whereas uniform WRRR and minimax have no hyper-parameters.}
\label{tab:stage1}
\centering\small
% AUTO-GENERATED by paper/tools/export_results.py
\begin{tabular}{lcccc}
\toprule
Method & Final rank & Test worst (mean $\pm$ sd) & Test mean & Cal.\ worst \\
\midrule
Base model ($W=0$) & 0 & 1.000 $\pm$ 0.000 & 1.000 & 1.000 \\
TA + SVD & 4 & 1.445 $\pm$ 0.472 & 0.770 & 1.400 \\
KnOTS-TA + trunc.$^\dagger$ & 4 & 1.445 $\pm$ 0.472 & 0.770 & 1.400 \\
KnOTS-TIES + trunc.$^\dagger$ & 4 & 1.077 $\pm$ 0.282 & 0.770 & 0.996 \\
CtM-like$^\dagger$ & 4 & 1.585 $\pm$ 0.421 & 0.835 & 1.543 \\
RegMean (full rank)$^\ddagger$ & 12 & 1.341 $\pm$ 0.506 & 0.701 & 0.948 \\
RegMean + Eucl.\ trunc. & 4 & 1.353 $\pm$ 0.414 & 0.734 & 1.012 \\
RegMean + whitened trunc. & 4 & 1.339 $\pm$ 0.401 & 0.739 & 1.012 \\
Uniform WRRR (rel.) & 4 & 0.685 $\pm$ 0.060 & 0.588 & 0.601 \\
Minimax (rel.) & 4 & 0.857 $\pm$ 0.179 & 0.627 & 0.472 \\
Factor averaging$^\S$ & 4 & 1.253 $\pm$ 0.350 & 0.747 & 1.204 \\
\bottomrule
\end{tabular}

\end{table}

\section{Correction history}
\label{app:corrections}

Summaries written before this draft contained the following discrepancies with the raw files. The
original summaries are preserved in the repository (\texttt{CORRECTIONS.md}).
\begin{enumerate}
\item The stage-1 minimax mean test worst-task error was summarised as ``0.858''. The raw value is
\SOneMinimaxRawMean{} (\texttt{per\_method.minimax\_rel.test\_worst\_rel.mean}), so the correct three-decimal
value is \SOneMinimaxTestWorstMean. The earlier value came from rounding a four-decimal display
twice.
\item The largest Monte Carlo deviation was summarised as ``3.03''. The raw maximum is
\STwoMCWorstZ; the earlier value was read from the last log line only.
\item The triples 0.713/0.802/1.057 and 0.726/0.768/1.081 were placed side by side in a summary.
They are the same minimax solutions evaluated on the test sample and on the population objective,
respectively (Table~\ref{tab:orig}). They are not a comparison of two methods.
\item An earlier interpretation said that the minimax solution failed because the weakest task
received a small weight. The oracle Pop/Pop solution also gives that task a small weight
($\lambda=\STwoPPLambdaTwoTA$ and $\STwoPPLambdaTwoTC$ for teachers 0 and 2). The small weight
itself is therefore not the failure.
\end{enumerate}

\section{Planned experiments (NOT RUN)}
\label{app:planned}

None of the following was executed. They are listed so that the absence of results is explicit.
\begin{enumerate}
\item \textbf{Real-adapter pilot} (BLOCKED). The plan is to merge \AdmNAdapters{} public Flan-T5-base GLUE
LoRA adapters (rank \AdmLoraR, query/value projections) and to report worst-task and mean task metrics
at a matched final rank and matched calibration access. It is blocked for three reasons:
\begin{itemize}
\item the adapter repositories contain no license or usage terms;
\item the base-model revision used for training is not recorded (the adapter configuration has
\texttt{revision: null});
\item the training and checkpoint-selection split is undocumented, while the publisher evaluates
on the GLUE validation split.
\end{itemize}
The metric plan must also separate the publisher's exact-match accuracy from the official GLUE
metrics (Section~\ref{sec:discussion}). The publisher reports the CoLA adapter at the same CoLA
score as the base model (\AdmCoLAScore), so accuracy normalised by the adapter's own score carries no
information for that task.
\item \textbf{Absolute versus relative errors}, to separate the effect of relative normalisation.
\item \textbf{Calibration-size dependence} (not approved in this phase: no $n$, rank or
task-count sweeps).
\item \textbf{Regularised worst-task variants}, for example shrinking $\lambda$ towards uniform.
These are not implemented; any such variant must be pre-registered and selected on development
data.
\item \textbf{Numerical evaluation of Corollary~\ref{cor:fixed} and Lemma~\ref{lem:denom}} on the
recorded instances. This is planned for the next revision of this draft.
\end{enumerate}

\section{Reproducibility details}
\label{app:repro}

\begin{itemize}
\item \emph{Stage 1.} Command:
\texttt{PYTHONPATH=src python3 -m lowrank\_merge.run\_cpu --config configs/cpu\_synthetic.json}.
Wall time \SOneWall~s.
\item \emph{Stage 2.} Registered before execution at commit \texttt{\STwoCommit}. Command:
\texttt{PYTHONPATH=src python3 -m lowrank\_merge.run\_diag --config
configs/p2\_stage2\_moment\_normalizer\_diag.json}.
\begin{itemize}
\item \STwoCellsOK{} of \STwoCellsTotal{} cells completed.
\item CPU time \STwoCPU~s under a hard limit of \STwoCPUCap~s; wall time \STwoWall~s; peak RSS
\STwoRSS~MiB.
\item One worker and one compute thread. No package was installed and nothing was downloaded.
\end{itemize}
\item \emph{Numbers and tables.} Every number in the text is emitted by
\texttt{paper/tools/export\_results.py}. The file \texttt{paper/generated/provenance.tsv} maps each
number to its source file and field.
\end{itemize}
